% Options for packages loaded elsewhere
% Options for packages loaded elsewhere
\PassOptionsToPackage{unicode}{hyperref}
\PassOptionsToPackage{hyphens}{url}
\PassOptionsToPackage{dvipsnames,svgnames,x11names}{xcolor}
%
\documentclass[
  russian,
]{article}
\usepackage{xcolor}
\usepackage[top=20mm,left=20mm,right=20mm,bottom=20mm]{geometry}
\usepackage{amsmath,amssymb}
\setcounter{secnumdepth}{5}
\usepackage{iftex}
\ifPDFTeX
  \usepackage[T1]{fontenc}
  \usepackage[utf8]{inputenc}
  \usepackage{textcomp} % provide euro and other symbols
\else % if luatex or xetex
  \usepackage{unicode-math} % this also loads fontspec
  \defaultfontfeatures{Scale=MatchLowercase}
  \defaultfontfeatures[\rmfamily]{Ligatures=TeX,Scale=1}
\fi
\usepackage{lmodern}
\ifPDFTeX\else
  % xetex/luatex font selection
  \setmainfont[]{Times New Roman}
  \setsansfont[]{Arial}
  \setmonofont[]{Courier New}
\fi
% Use upquote if available, for straight quotes in verbatim environments
\IfFileExists{upquote.sty}{\usepackage{upquote}}{}
\IfFileExists{microtype.sty}{% use microtype if available
  \usepackage[]{microtype}
  \UseMicrotypeSet[protrusion]{basicmath} % disable protrusion for tt fonts
}{}
\makeatletter
\@ifundefined{KOMAClassName}{% if non-KOMA class
  \IfFileExists{parskip.sty}{%
    \usepackage{parskip}
  }{% else
    \setlength{\parindent}{0pt}
    \setlength{\parskip}{6pt plus 2pt minus 1pt}}
}{% if KOMA class
  \KOMAoptions{parskip=half}}
\makeatother
% Make \paragraph and \subparagraph free-standing
\makeatletter
\ifx\paragraph\undefined\else
  \let\oldparagraph\paragraph
  \renewcommand{\paragraph}{
    \@ifstar
      \xxxParagraphStar
      \xxxParagraphNoStar
  }
  \newcommand{\xxxParagraphStar}[1]{\oldparagraph*{#1}\mbox{}}
  \newcommand{\xxxParagraphNoStar}[1]{\oldparagraph{#1}\mbox{}}
\fi
\ifx\subparagraph\undefined\else
  \let\oldsubparagraph\subparagraph
  \renewcommand{\subparagraph}{
    \@ifstar
      \xxxSubParagraphStar
      \xxxSubParagraphNoStar
  }
  \newcommand{\xxxSubParagraphStar}[1]{\oldsubparagraph*{#1}\mbox{}}
  \newcommand{\xxxSubParagraphNoStar}[1]{\oldsubparagraph{#1}\mbox{}}
\fi
\makeatother


\usepackage{longtable,booktabs,array}
\newcounter{none} % for unnumbered tables
\usepackage{calc} % for calculating minipage widths
% Correct order of tables after \paragraph or \subparagraph
\usepackage{etoolbox}
\makeatletter
\patchcmd\longtable{\par}{\if@noskipsec\mbox{}\fi\par}{}{}
\makeatother
% Allow footnotes in longtable head/foot
\IfFileExists{footnotehyper.sty}{\usepackage{footnotehyper}}{\usepackage{footnote}}
\makesavenoteenv{longtable}
\usepackage{graphicx}
\makeatletter
\newsavebox\pandoc@box
\newcommand*\pandocbounded[1]{% scales image to fit in text height/width
  \sbox\pandoc@box{#1}%
  \Gscale@div\@tempa{\textheight}{\dimexpr\ht\pandoc@box+\dp\pandoc@box\relax}%
  \Gscale@div\@tempb{\linewidth}{\wd\pandoc@box}%
  \ifdim\@tempb\p@<\@tempa\p@\let\@tempa\@tempb\fi% select the smaller of both
  \ifdim\@tempa\p@<\p@\scalebox{\@tempa}{\usebox\pandoc@box}%
  \else\usebox{\pandoc@box}%
  \fi%
}
% Set default figure placement to htbp
\def\fps@figure{htbp}
\makeatother



\ifLuaTeX
\usepackage[bidi=basic,shorthands=off,provide=*]{babel}
\else
\usepackage[bidi=default,shorthands=off,provide=*]{babel}
\fi
\ifPDFTeX
\else
\babelfont{rm}[]{Times New Roman}
\fi
\ifLuaTeX
  \usepackage{selnolig} % disable illegal ligatures
\fi


\setlength{\emergencystretch}{3em} % prevent overfull lines

\providecommand{\tightlist}{%
  \setlength{\itemsep}{0pt}\setlength{\parskip}{0pt}}



 


\makeatletter
\@ifpackageloaded{caption}{}{\usepackage{caption}}
\AtBeginDocument{%
\ifdefined\contentsname
  \renewcommand*\contentsname{Содержание}
\else
  \newcommand\contentsname{Содержание}
\fi
\ifdefined\listfigurename
  \renewcommand*\listfigurename{Список Иллюстраций}
\else
  \newcommand\listfigurename{Список Иллюстраций}
\fi
\ifdefined\listtablename
  \renewcommand*\listtablename{Список Таблиц}
\else
  \newcommand\listtablename{Список Таблиц}
\fi
\ifdefined\figurename
  \renewcommand*\figurename{Рисунок}
\else
  \newcommand\figurename{Рисунок}
\fi
\ifdefined\tablename
  \renewcommand*\tablename{Таблица}
\else
  \newcommand\tablename{Таблица}
\fi
}
\@ifpackageloaded{float}{}{\usepackage{float}}
\floatstyle{ruled}
\@ifundefined{c@chapter}{\newfloat{codelisting}{h}{lop}}{\newfloat{codelisting}{h}{lop}[chapter]}
\floatname{codelisting}{Список}
\newcommand*\listoflistings{\listof{codelisting}{Список Каталогов}}
\makeatother
\makeatletter
\makeatother
\makeatletter
\@ifpackageloaded{caption}{}{\usepackage{caption}}
\@ifpackageloaded{subcaption}{}{\usepackage{subcaption}}
\makeatother
\usepackage{bookmark}
\IfFileExists{xurl.sty}{\usepackage{xurl}}{} % add URL line breaks if available
\urlstyle{same}
\hypersetup{
  pdftitle={Задачи по Эконометрике-2: Прогнозирование для logit/probit-моделей},
  pdfauthor={Н.В. Артамонов, А. Н. Лата (МГИМО МИД России)},
  pdflang={ru},
  colorlinks=true,
  linkcolor={blue},
  filecolor={Maroon},
  citecolor={Blue},
  urlcolor={Blue},
  pdfcreator={LaTeX via pandoc}}


\title{Задачи по Эконометрике-2: Прогнозирование для
logit/probit-моделей}
\author{Н.В. Артамонов, А. Н. Лата (МГИМО МИД России)}
\date{}
\begin{document}
\maketitle

\renewcommand*\contentsname{Содержание}
{
\setcounter{tocdepth}{3}
\tableofcontents
}

\section{labour force equation}\label{labour-force-equation}

Для датасета \texttt{TableF5-1} рассмотрим регрессию \textbf{LFP на WA,
log(FAMINC), WE, KL6, K618, CIT, UN} трёх спецификаций: - LPM - logit -
probit

Результаты подгонки:

{\def\LTcaptype{none} % do not increment counter
\begin{longtable}[]{@{}llll@{}}
\toprule\noalign{}
& LPM & logit & probit \\
\midrule\noalign{}
\endhead
\bottomrule\noalign{}
\endlastfoot
Dep. Variable & LFP & LFP & LFP \\
WA & -0.013*** & -0.063*** & -0.038*** \\
& (0.003) & (0.013) & (0.008) \\
np.log(FAMINC) & 0.075** & 0.341** & 0.205** \\
& (0.037) & (0.172) & (0.104) \\
WE & 0.038*** & 0.179*** & 0.108*** \\
& (0.008) & (0.040) & (0.024) \\
KL6 & -0.302*** & -1.443*** & -0.868*** \\
& (0.036) & (0.194) & (0.112) \\
K618 & -0.018 & -0.095 & -0.057 \\
& (0.014) & (0.067) & (0.040) \\
CIT & -0.048 & -0.214 & -0.126 \\
& (0.038) & (0.176) & (0.107) \\
UN & -0.004 & -0.017 & -0.011 \\
& (0.006) & (0.026) & (0.016) \\
Intercept & 0.079 & -1.856 & -1.108 \\
& (0.356) & (1.679) & (1.014) \\
R-squared & 0.130 & & \\
R-squared Adj. & 0.122 & & \\
\end{longtable}
}

\emph{Signif. codes:} \texttt{***} \(p<0.01\), \texttt{**} \(p<0.05\),
\texttt{*} \(p<0.1\)

Рассмотрим несколько человек с характеристиками. Постройте прогноз для
каждого. \textbf{Ответ округлите до 3-х десятичных знаков.}

\textbf{Характеристики:}

{\def\LTcaptype{none} % do not increment counter
\begin{longtable}[]{@{}rrrrrrr@{}}
\toprule\noalign{}
WA & FAMINC & WE & KL6 & K618 & CIT & UN \\
\midrule\noalign{}
\endhead
\bottomrule\noalign{}
\endlastfoot
35 & 12500 & 15 & 2 & 0 & 1 & 5 \\
40 & 9800 & 12 & 1 & 2 & 0 & 7.5 \\
42 & 67800 & 14 & 2 & 1 & 1 & 3 \\
\end{longtable}
}

\textbf{Ответ:}

{\def\LTcaptype{none} % do not increment counter
\begin{longtable}[]{@{}rrrr@{}}
\toprule\noalign{}
Человек & Прогноз LPM & Прогноз logit & Прогноз probit \\
\midrule\noalign{}
\endhead
\bottomrule\noalign{}
\endlastfoot
1 & 0.221 & 0.209 & 0.214 \\
2 & 0.329 & 0.3 & 0.307 \\
3 & 0.207 & 0.192 & 0.196 \\
\end{longtable}
}

\section{approve equation}\label{approve-equation}

Для датасета \texttt{loanapp} рассмотрим регрессию \textbf{approve на
appinc/100, mortno, unem, dep, male} трёх спецификаций: - LPM - logit -
probit

Результаты подгонки:

{\def\LTcaptype{none} % do not increment counter
\begin{longtable}[]{@{}llll@{}}
\toprule\noalign{}
& LPM & logit & probit \\
\midrule\noalign{}
\endhead
\bottomrule\noalign{}
\endlastfoot
Dep. Variable & approve & approve & approve \\
I(appinc / 100) & -0.014* & -0.106 & -0.056 \\
& (0.009) & (0.067) & (0.038) \\
mortno & 0.079*** & 0.817*** & 0.422*** \\
& (0.016) & (0.170) & (0.086) \\
unem & -0.008** & -0.065** & -0.036** \\
& (0.003) & (0.029) & (0.016) \\
dep & -0.012* & -0.106* & -0.055* \\
& (0.007) & (0.061) & (0.033) \\
male & 0.019 & 0.173 & 0.093 \\
& (0.019) & (0.176) & (0.094) \\
Intercept & 0.886*** & 2.032*** & 1.198*** \\
& (0.022) & (0.193) & (0.105) \\
R-squared & 0.017 & & \\
R-squared Adj. & 0.014 & & \\
\end{longtable}
}

\emph{Signif. codes:} \texttt{***} \(p<0.01\), \texttt{**} \(p<0.05\),
\texttt{*} \(p<0.1\)

Рассмотрим несколько человек с характеристиками. Постройте прогноз для
каждого. \textbf{Ответ округлите до 3-х десятичных знаков.}

\textbf{Характеристики:}

{\def\LTcaptype{none} % do not increment counter
\begin{longtable}[]{@{}rrrrr@{}}
\toprule\noalign{}
appinc & mortno & unem & dep & male \\
\midrule\noalign{}
\endhead
\bottomrule\noalign{}
\endlastfoot
120 & 1 & 1.8 & 0 & 1 \\
48 & 1 & 3.2 & 0 & 0 \\
82 & 0 & 3.9 & 1 & 1 \\
\end{longtable}
}

\textbf{Ответ:}

{\def\LTcaptype{none} % do not increment counter
\begin{longtable}[]{@{}rrrr@{}}
\toprule\noalign{}
Человек & Прогноз LPM & Прогноз logit & Прогноз probit \\
\midrule\noalign{}
\endhead
\bottomrule\noalign{}
\endlastfoot
1 & 0.953 & 0.941 & 0.943 \\
2 & 0.933 & 0.93 & 0.93 \\
3 & 0.851 & 0.853 & 0.853 \\
\end{longtable}
}

\section{swiss labour force equation}\label{swiss-labour-force-equation}

Для датасета \texttt{SwissLabor} рассмотрим регрессию
\textbf{participation на income, income\^{}2, age, age\^{}2, youngkids,
oldkids} трёх спецификаций: - LPM - logit - probit

Результаты подгонки:

{\def\LTcaptype{none} % do not increment counter
\begin{longtable}[]{@{}llll@{}}
\toprule\noalign{}
& LPM & logit & probit \\
\midrule\noalign{}
\endhead
\bottomrule\noalign{}
\endlastfoot
Dep. Variable & participation & participation & participation \\
income & 0.269 & 3.146 & 1.973 \\
& (0.654) & (3.400) & (2.090) \\
I(income ** 2) & -0.025 & -0.211 & -0.131 \\
& (0.031) & (0.161) & (0.099) \\
age & 0.790*** & 3.875*** & 2.347*** \\
& (0.131) & (0.669) & (0.397) \\
I(age ** 2) & -0.111*** & -0.544*** & -0.329*** \\
& (0.016) & (0.083) & (0.049) \\
youngkids & -0.227*** & -1.065*** & -0.637*** \\
& (0.032) & (0.166) & (0.096) \\
oldkids & -0.052*** & -0.250*** & -0.151*** \\
& (0.017) & (0.083) & (0.050) \\
Intercept & -0.716 & -15.282 & -9.643 \\
& (3.488) & (17.960) & (11.066) \\
R-squared & 0.155 & & \\
R-squared Adj. & 0.149 & & \\
\end{longtable}
}

\emph{Signif. codes:} \texttt{***} \(p<0.01\), \texttt{**} \(p<0.05\),
\texttt{*} \(p<0.1\)

Рассмотрим несколько человек с характеристиками. Постройте прогноз для
каждого. \textbf{Ответ округлите до 3-х десятичных знаков.}

\textbf{Характеристики:}

{\def\LTcaptype{none} % do not increment counter
\begin{longtable}[]{@{}rrrr@{}}
\toprule\noalign{}
income & age & youngkids & oldkids \\
\midrule\noalign{}
\endhead
\bottomrule\noalign{}
\endlastfoot
11.367 & 2.5 & 0 & 0 \\
9.217 & 3.7 & 2 & 0 \\
10.686 & 4.2 & 2 & 1 \\
\end{longtable}
}

\textbf{Ответ:}

{\def\LTcaptype{none} % do not increment counter
\begin{longtable}[]{@{}rrrr@{}}
\toprule\noalign{}
Человек & Прогноз LPM & Прогноз logit & Прогноз probit \\
\midrule\noalign{}
\endhead
\bottomrule\noalign{}
\endlastfoot
1 & 0.419 & 0.367 & 0.371 \\
2 & 0.606 & 0.626 & 0.624 \\
3 & 0.18 & 0.182 & 0.186 \\
\end{longtable}
}




\end{document}
